% GENERATED FILE. Edit canonical lessons, then run npm run generate:books.
% canonical-source-hash: fnv1a64:2f7ce5c01ba78082
% canonical-lessons: MR-C449-nisarga, MR-C449-samvardhan, MR-C449-vrksaropan, MR-C449-jangaltod, MR-C449-namses-hone, MR-R449-review-443-449, MR-R449-review-final-449

\chapter{Keeping Nature Green}
\label{ch:mr-keeping-nature-green}
\hlchaptermodality{449}

\begin{chapteropening}
\textbf{By the end of this chapter:} \emph{I can talk about nature, conservation, tree planting and deforestation, and say a species died out.}

Complete the last lesson of chapter 449: I can talk about nature, conservation, tree planting and deforestation, and say a species died out.
\end{chapteropening}

\section[\texorpdfstring{\mr{निसर्ग} (nisarga)}{निसर्ग (nisarga)}]{\mr{निसर्ग} (nisarga) — nature}

\label{lesson:MR-C449-nisarga}



\begin{quote}
Before the new one: say the Marathi for to rot, then the Marathi for to decay.
\end{quote}

\subsection*{You'll want to know: \mr{निसर्ग}}
\textbf{\mr{निसर्ग}} — \emph{nisarga} — \textquotedblleft{}nature\textquotedblright{}.

\begin{quote}
\mr{मला} \mr{निसर्गात} \mr{फिरायला} \mr{आवडतं}.
\end{quote}

\begin{quote}
\emph{malā nisargāt phirāylā āvaḍtaṁ.}
\end{quote}

\begin{quote}
I like walking out in nature.
\end{quote}

Masculine.

\mr{निसर्गरम्य}, beautiful with nature, is a common word for scenery.

\begin{cousinweb}
\textbf{\mr{निसर्ग}} (\emph{nisarga}) is a tatsama, taken from Sanskrit \emph{nisarga}, \textquotedblleft{}giving forth; the natural state\textquotedblright{} (Monier-Williams).
\end{cousinweb}

\begin{grammarlens}[title={What you've built}]
One more word for getting around.
\end{grammarlens}

\subsection*{Your turn}
\begin{itemize}
\raggedright
  \item \emph{Say it:} \emph{nisarga}
  \item \emph{Say it:} \emph{nisarga}, once more
  \item \emph{Say it:} \emph{malā nisargāt phirāylā āvaḍtaṁ.} — I like walking out in nature.
  \item \emph{Recall:} say the Marathi for to rot, then the Marathi for to decay, then say \emph{nisarga} again
\end{itemize}

\begin{tcolorbox}[breakable,colback=teal!4,colframe=teal!35!black,title={Before you move on}]
What does \mr{निसर्ग} mean? (\textquotedblleft{}Nature\textquotedblright{}.) Say it once more.
\end{tcolorbox}
\section[\texorpdfstring{\mr{संवर्धन} (saṁvardhan)}{संवर्धन (saṁvardhan)}]{\mr{संवर्धन} (saṁvardhan) — conservation}

\label{lesson:MR-C449-samvardhan}



\begin{quote}
Before the new one: say the Marathi for to decay, then the Marathi for nature.
\end{quote}

\subsection*{You'll want to know: \mr{संवर्धन}}
\textbf{\mr{संवर्धन}} — \emph{saṁvardhan} — \textquotedblleft{}conservation\textquotedblright{}.

\begin{quote}
\mr{वाघांच्या} \mr{संवर्धनासाठी} \mr{हे} \mr{अभयारण्य} \mr{आहे}.
\end{quote}

\begin{quote}
\emph{vāghāñcyā saṁvardhanāsāṭhī he abhayāraṇya āhe.}
\end{quote}

\begin{quote}
This sanctuary is for the conservation of tigers.
\end{quote}

Neuter.

\mr{संवर्धन} \mr{करणे}, to conserve and help grow.

\begin{cousinweb}
\textbf{\mr{संवर्धन}} (\emph{saṁvardhan}) is a tatsama, taken from Sanskrit \emph{saṁvardhana}, \textquotedblleft{}growing up, promoting, rearing\textquotedblright{} (Monier-Williams).
\end{cousinweb}

\begin{grammarlens}[title={What you've built}]
One more word for getting around.
\end{grammarlens}

\subsection*{Your turn}
\begin{itemize}
\raggedright
  \item \emph{Say it:} \emph{saṁvardhan}
  \item \emph{Say it:} \emph{saṁvardhan}, once more
  \item \emph{Say it:} \emph{vāghāñcyā saṁvardhanāsāṭhī he abhayāraṇya āhe.} — This sanctuary is for the conservation of tigers.
  \item \emph{Recall:} say the Marathi for to decay, then the Marathi for nature, then say \emph{saṁvardhan} again
\end{itemize}

\begin{tcolorbox}[breakable,colback=teal!4,colframe=teal!35!black,title={Before you move on}]
What does \mr{संवर्धन} mean? (\textquotedblleft{}Conservation\textquotedblright{}.) Say it once more.
\end{tcolorbox}
\section[\texorpdfstring{\mr{वृक्षारोपण} (vṛkṣāropaṇ)}{वृक्षारोपण (vṛkṣāropaṇ)}]{\mr{वृक्षारोपण} (vṛkṣāropaṇ) — tree planting}

\label{lesson:MR-C449-vrksaropan}



\begin{quote}
Before the new one: say the Marathi for nature, then the Marathi for conservation.
\end{quote}

\subsection*{You'll want to know: \mr{वृक्षारोपण}}
\textbf{\mr{वृक्षारोपण}} — \emph{vṛkṣāropaṇ} — \textquotedblleft{}tree planting\textquotedblright{}.

\begin{quote}
\mr{शाळेने} \mr{वृक्षारोपणाचा} \mr{कार्यक्रम} \mr{घेतला}.
\end{quote}

\begin{quote}
\emph{śāḷene vṛkṣāropaṇācā kāryakram ghetlā.}
\end{quote}

\begin{quote}
The school held a tree planting event.
\end{quote}

Neuter.

\mr{वृक्ष}, a tree, and \mr{आरोपण}, planting.

\begin{cousinweb}
\textbf{\mr{वृक्षारोपण}} (\emph{vṛkṣāropaṇ}) is a tatsama, taken from Sanskrit \emph{vṛkṣāropaṇa}, \textquotedblleft{}planting trees\textquotedblright{} (Monier-Williams).
\end{cousinweb}

\begin{grammarlens}[title={What you've built}]
One more word for getting around.
\end{grammarlens}

\subsection*{Your turn}
\begin{itemize}
\raggedright
  \item \emph{Say it:} \emph{vṛkṣāropaṇ}
  \item \emph{Say it:} \emph{vṛkṣāropaṇ}, once more
  \item \emph{Say it:} \emph{śāḷene vṛkṣāropaṇācā kāryakram ghetlā.} — The school held a tree planting event.
  \item \emph{Recall:} say the Marathi for nature, then the Marathi for conservation, then say \emph{vṛkṣāropaṇ} again
\end{itemize}

\begin{tcolorbox}[breakable,colback=teal!4,colframe=teal!35!black,title={Before you move on}]
What does \mr{वृक्षारोपण} mean? (\textquotedblleft{}Tree planting\textquotedblright{}.) Say it once more.
\end{tcolorbox}
\section[\texorpdfstring{\mr{जंगलतोड} (jaṅgaltoḍ)}{जंगलतोड (jaṅgaltoḍ)}]{\mr{जंगलतोड} (jaṅgaltoḍ) — deforestation}

\label{lesson:MR-C449-jangaltod}



\begin{quote}
Before the new one: say the Marathi for conservation, then the Marathi for tree planting.
\end{quote}

\subsection*{You'll want to know: \mr{जंगलतोड}}
\textbf{\mr{जंगलतोड}} — \emph{jaṅgaltoḍ} — \textquotedblleft{}deforestation\textquotedblright{}.

\begin{quote}
\mr{जंगलतोडीमुळे} \mr{पाऊस} \mr{कमी} \mr{झाला}.
\end{quote}

\begin{quote}
\emph{jaṅgaltoḍīmuḷe pāūs kamī jhālā.}
\end{quote}

\begin{quote}
The rain has decreased because of deforestation.
\end{quote}

Feminine.

Before an ending \mr{जंगलतोडी}-: \mr{जंगलतोडीमुळे}.

\begin{cousinweb}
\textbf{\mr{जंगलतोड}} (\emph{jaṅgaltoḍ}) is built from \textbf{\mr{जंगल}}, \textquotedblleft{}a forest\textquotedblright{}, from the family of Sanskrit \emph{jāṅgala}, \textquotedblleft{}arid land\textquotedblright{}, where Hindi has \emph{jaṅgal}, \textquotedblleft{}forest\textquotedblright{} (CDIAL 5177), and \emph{toḍ}, \textquotedblleft{}cutting\textquotedblright{}, from \textbf{\mr{तोडणे}}.
\end{cousinweb}

\begin{grammarlens}[title={What you've built}]
One more word for getting around.
\end{grammarlens}

\subsection*{Your turn}
\begin{itemize}
\raggedright
  \item \emph{Say it:} \emph{jaṅgaltoḍ}
  \item \emph{Say it:} \emph{jaṅgaltoḍ}, once more
  \item \emph{Say it:} \emph{jaṅgaltoḍīmuḷe pāūs kamī jhālā.} — The rain has decreased because of deforestation.
  \item \emph{Recall:} say the Marathi for conservation, then the Marathi for tree planting, then say \emph{jaṅgaltoḍ} again
\end{itemize}

\begin{tcolorbox}[breakable,colback=teal!4,colframe=teal!35!black,title={Before you move on}]
What does \mr{जंगलतोड} mean? (\textquotedblleft{}Deforestation\textquotedblright{}.) Say it once more.
\end{tcolorbox}
\section[\texorpdfstring{\mr{नामशेष} \mr{होणे} (nāmśeṣ hoṇe)}{नामशेष होणे (nāmśeṣ hoṇe)}]{\mr{नामशेष} \mr{होणे} (nāmśeṣ hoṇe) — to die out}

\label{lesson:MR-C449-namses-hone}



\begin{quote}
Before the new one: say the Marathi for tree planting, then the Marathi for deforestation.
\end{quote}

\subsection*{You'll want to know: \mr{नामशेष} \mr{होणे}}
\textbf{\mr{नामशेष} \mr{होणे}} — \emph{nāmśeṣ hoṇe} — \textquotedblleft{}to die out\textquotedblright{}.

\begin{quote}
\mr{हा} \mr{पक्षी} \mr{आता} \mr{नामशेष} \mr{झाला} \mr{आहे}.
\end{quote}

\begin{quote}
\emph{hā pakṣī ātā nāmśeṣ jhālā āhe.}
\end{quote}

\begin{quote}
This bird has now died out.
\end{quote}

Past (he …): \mr{तो} \mr{नामशेष} \mr{झाला}, \emph{to nāmśeṣ jhālā}.

\mr{नाम}, the name, and \mr{शेष}, what is left: only the name remains.

\begin{cousinweb}
\textbf{\mr{नामशेष} \mr{होणे}} (\emph{nāmśeṣ hoṇe}) joins \emph{nāmaśeṣa}, \textquotedblleft{}of whom only the name is left\textquotedblright{}, built from Sanskrit \emph{nāman}, \textquotedblleft{}name\textquotedblright{}, and \emph{śeṣa}, \textquotedblleft{}remainder\textquotedblright{} (Monier-Williams), to \textbf{\mr{होणे}}.
\end{cousinweb}

\begin{grammarlens}[title={What you've built}]
One more word for getting around.
\end{grammarlens}

\subsection*{Your turn}
\begin{itemize}
\raggedright
  \item \emph{Say it:} \emph{nāmśeṣ hoṇe}
  \item \emph{Say it:} \emph{nāmśeṣ hoṇe}, once more
  \item \emph{Say it:} \emph{hā pakṣī ātā nāmśeṣ jhālā āhe.} — This bird has now died out.
  \item \emph{Recall:} say the Marathi for tree planting, then the Marathi for deforestation, then say \emph{nāmśeṣ hoṇe} again
\end{itemize}

\begin{tcolorbox}[breakable,colback=teal!4,colframe=teal!35!black,title={Before you move on}]
What does \mr{नामशेष} \mr{होणे} mean? (\textquotedblleft{}To die out\textquotedblright{}.) Say it once more.
\end{tcolorbox}
\section[(dialogue)]{Looking Back on the Wild}

\label{lesson:MR-R449-review-443-449}



\begin{quote}
Say the Marathi for a stream and a canal.
\end{quote}

\subsection*{Your turn}
Say these aloud:

\begin{itemize}
\raggedright
  \item a stream, a canal, a dam, a large lake, to gurgle
  \item a puddle, silt, muddy, to seep, to overflow
  \item a sprout, to take root, to bloom, to blossom, to wither
  \item wildlife, a sanctuary, a leopard, to roar, to graze
  \item a heron, a black kite, a vulture, to chirp, migration
  \item waste water, a stench, polluted, to rot, to decay
  \item nature, conservation, tree planting, deforestation, to die out
\end{itemize}

\begin{tcolorbox}[breakable,colback=teal!4,colframe=teal!35!black,title={Before you move on}]
A stream? (\textbf{\mr{ओढा}}, \emph{oḍhā}.) A canal? (\textbf{\mr{कालवा}}, \emph{kālvā}.) A dam? (\textbf{\mr{धरण}}, \emph{dharaṇ}.) A large lake? (\textbf{\mr{सरोवर}}, \emph{sarovar}.) To gurgle? (\textbf{\mr{खळखळणे}}, \emph{khaḷkhaḷṇe}.) A puddle? (\textbf{\mr{डबके}}, \emph{ḍabke}.) Silt? (\textbf{\mr{गाळ}}, \emph{gāḷ}.) Muddy? (\textbf{\mr{गढूळ}}, \emph{gaḍhūḷ}.) To seep? (\textbf{\mr{झिरपणे}}, \emph{jhirapṇe}.) To overflow? (\textbf{\mr{ओसंडणे}}, \emph{osaṇḍṇe}.) A sprout? (\textbf{\mr{अंकुर}}, \emph{aṅkur}.) To take root? (\textbf{\mr{रुजणे}}, \emph{rujṇe}.) To bloom? (\textbf{\mr{फुलणे}}, \emph{phulṇe}.) To blossom? (\textbf{\mr{बहरणे}}, \emph{baharṇe}.) To wither? (\textbf{\mr{कोमेजणे}}, \emph{komejṇe}.) Wildlife? (\textbf{\mr{वन्यजीव}}, \emph{vanyajīv}.) A sanctuary? (\textbf{\mr{अभयारण्य}}, \emph{abhayāraṇya}.) A leopard? (\textbf{\mr{बिबट्या}}, \emph{bibaṭyā}.) To roar? (\textbf{\mr{डरकाळी} \mr{फोडणे}}, \emph{ḍarkāḷī phoḍṇe}.) To graze? (\textbf{\mr{चरणे}}, \emph{carṇe}.) A heron? (\textbf{\mr{बगळा}}, \emph{bagḷā}.) A black kite? (\textbf{\mr{घार}}, \emph{ghār}.) A vulture? (\textbf{\mr{गिधाड}}, \emph{gidhāḍ}.) To chirp? (\textbf{\mr{किलबिलणे}}, \emph{kilbilṇe}.) Migration? (\textbf{\mr{स्थलांतर}}, \emph{sthalāntar}.) Waste water? (\textbf{\mr{सांडपाणी}}, \emph{sāṇḍpāṇī}.) A stench? (\textbf{\mr{दुर्गंधी}}, \emph{durgandhī}.) Polluted? (\textbf{\mr{प्रदूषित}}, \emph{pradūṣit}.) To rot? (\textbf{\mr{कुजणे}}, \emph{kujṇe}.) To decay? (\textbf{\mr{सडणे}}, \emph{saḍṇe}.) Nature? (\textbf{\mr{निसर्ग}}, \emph{nisarga}.) Conservation? (\textbf{\mr{संवर्धन}}, \emph{saṁvardhan}.) Tree planting? (\textbf{\mr{वृक्षारोपण}}, \emph{vṛkṣāropaṇ}.) Deforestation? (\textbf{\mr{जंगलतोड}}, \emph{jaṅgaltoḍ}.) To die out? (\textbf{\mr{नामशेष} \mr{होणे}}, \emph{nāmśeṣ hoṇe}.)
\end{tcolorbox}
\section[(dialogue)]{Keeping Nature Green — once more}

\label{lesson:MR-R449-review-final-449}



\begin{quote}
Say the Marathi for nature and conservation.
\end{quote}

\subsection*{Your turn}
\begin{itemize}
\raggedright
  \item \emph{Say it:} nature, conservation, tree planting, deforestation, to die out
\end{itemize}

\begin{tcolorbox}[breakable,colback=teal!4,colframe=teal!35!black,title={Before you move on}]
Nature? (\textbf{\mr{निसर्ग}}, \emph{nisarga}.) Conservation? (\textbf{\mr{संवर्धन}}, \emph{saṁvardhan}.) Tree planting? (\textbf{\mr{वृक्षारोपण}}, \emph{vṛkṣāropaṇ}.) Deforestation? (\textbf{\mr{जंगलतोड}}, \emph{jaṅgaltoḍ}.) To die out? (\textbf{\mr{नामशेष} \mr{होणे}}, \emph{nāmśeṣ hoṇe}.)
\end{tcolorbox}
